\documentclass[12pt,a4paper]{article}

% =============================================================================
% PAQUETES
% =============================================================================
\usepackage[utf8]{inputenc}
\usepackage[T1]{fontenc}
\usepackage[spanish]{babel}
\usepackage{geometry}
\usepackage{graphicx}
\usepackage{xcolor}
\usepackage{booktabs}
\usepackage{longtable}
\usepackage{array}
\usepackage{multirow}
\usepackage{float}
\usepackage{caption}
\usepackage{subcaption}
\usepackage{amsmath}
\usepackage{amssymb}
\usepackage{enumitem}
\usepackage{fancyhdr}
\usepackage{titlesec}
\usepackage{setspace}
\usepackage{tikz}
\usepackage{listings}
\usepackage{pgfplots}
\pgfplotsset{compat=1.18}

% =============================================================================
% CONFIGURACIÓN DE PÁGINA
% =============================================================================
\geometry{left=2.5cm, right=2.5cm, top=2.5cm, bottom=2.5cm}

% =============================================================================
% COLORES
% =============================================================================
\definecolor{azuloscuro}{RGB}{0, 51, 102}
\definecolor{azulclaro}{RGB}{70, 130, 180}
\definecolor{rojoacento}{RGB}{178, 34, 34}
\definecolor{verdeok}{RGB}{0, 128, 0}
\definecolor{codebg}{RGB}{248,248,248}
\definecolor{codeframe}{RGB}{220,220,220}
\definecolor{codecomment}{RGB}{0,128,0}
\definecolor{codestring}{RGB}{163,21,21}
\definecolor{codekeyword}{RGB}{0,0,255}

\lstset{
    language=R,
    backgroundcolor=\color{codebg},
    frame=single,
    rulecolor=\color{codeframe},
    basicstyle=\ttfamily\small,
    keywordstyle=\color{codekeyword}\bfseries,
    commentstyle=\color{codecomment}\itshape,
    stringstyle=\color{codestring},
    numbers=left,
    numberstyle=\tiny\color{gray},
    numbersep=8pt,
    breaklines=true,
    breakatwhitespace=true,
    showstringspaces=false,
    tabsize=2,
    captionpos=b,
    xleftmargin=15pt,
    xrightmargin=5pt,
    aboveskip=10pt,
    belowskip=10pt
}

% =============================================================================
% ENCABEZADOS Y TÍTULOS
% =============================================================================
\pagestyle{fancy}
\fancyhf{}
\fancyhead[L]{\textcolor{azuloscuro}{\small Análisis de Datos en Ingeniería I}}
\fancyhead[R]{\textcolor{azuloscuro}{\small Guía Didáctica --- Código R y RLM}}
\fancyfoot[C]{\thepage}
\renewcommand{\headrulewidth}{0.5pt}
\renewcommand{\footrulewidth}{0.3pt}

\titleformat{\section}{\normalfont\Large\bfseries\color{azuloscuro}}{\thesection.}{0.5em}{}[\titlerule]
\titleformat{\subsection}{\normalfont\large\bfseries\color{azulclaro}}{\thesubsection}{0.5em}{}
\titleformat{\subsubsection}{\normalfont\normalsize\bfseries\color{azuloscuro}}{\thesubsubsection}{0.5em}{}

\setstretch{1.5}

% =============================================================================
% INICIO DEL DOCUMENTO
% =============================================================================
\begin{document}

% PORTADA
\begin{titlepage}
    \begin{tikzpicture}[remember picture, overlay]
        \fill[azuloscuro] (current page.north west) rectangle ([yshift=-3cm]current page.north east);
        \fill[azuloscuro] (current page.south west) rectangle ([yshift=2cm]current page.south east);
    \end{tikzpicture}
    \vspace*{4cm}
    \begin{center}
        {\Huge\bfseries\textcolor{azuloscuro}{GUÍA DIDÁCTICA}}\\[0.5cm]
        {\LARGE\bfseries\textcolor{azulclaro}{Análisis de Datos en Ingeniería I}}\\[1cm]
        \rule{0.8\textwidth}{1.5pt}\\[0.8cm]
        {\Large\textbf{Explicación Profunda del Código R}}\\[0.3cm]
        {\Large\textbf{y Resultados Reales del Proyecto}}\\[1.5cm]
        \rule{0.6\textwidth}{0.5pt}\\[1cm]
        {\large\textbf{Regresión Lineal Múltiple}}\\[0.3cm]
        {\large\textbf{Predicción de Precios de Laptops}}\\[1.5cm]
        {\large\textbf{Integrantes:}}\\[0.5cm]
        \begin{tabular}{rl}
            Dubal Aguilar Torres & (200207378) \\
            Alejandro Chaves Ramos & (200116116) \\
            Juan Cáceres Figueroa & (200079919) \\
            Miguel Carrizosa Morales & (200208123) \\
        \end{tabular}
        \vfill
        {\large\textbf{Docente:}}\\[0.3cm]
        PhD. Luis Ángel Anillo Arrieta\\[1cm]
        {\large Ingeniería de Sistemas}\\[0.3cm]
        {\large Universidad del Norte}\\[0.3cm]
        {\large Barranquilla, Colombia}\\[0.5cm]
        {\large 25 de mayo de 2026}
    \end{center}
\end{titlepage}

\newpage
\tableofcontents
\newpage

% =============================================================================
% 1. INTRODUCCIÓN
% =============================================================================
\section{Introducción}

Esta guía tiene un objetivo puramente didáctico: que entiendas a profundidad qué hace cada línea del código R del proyecto, por qué se hace, y qué significa estadísticamente. Incluye los \textbf{resultados reales} ejecutados con el dataset de 1,303 laptops.

\subsection{Dataset utilizado}

\begin{itemize}
    \item \textbf{Fuente:} Laptop Price Prediction Dataset (Kaggle)
    \item \textbf{Registros totales:} 1,303 laptops
    \item \textbf{Registros tras limpieza:} 1,303 (no hubo datos faltantes en las variables del modelo)
    \item \textbf{Variables:} 13 columnas originales; se crearon variables adicionales tras limpieza
\end{itemize}

\subsection{Modelo ejecutado}

El modelo ajustado es:
\[
\text{Price\_euros}_i = \beta_0 + \beta_1 \cdot \text{RAM}_i + \beta_2 \cdot \text{CPU\_GHz}_i + \beta_3 \cdot \text{Inches}_i + \beta_{\text{TypeName}} + \beta_{\text{OpSys}} + \varepsilon_i
\]

Con $N = 1{,}303$ observaciones y $k = 16$ predictores (incluyendo dummies).

% =============================================================================
% 2. LIBRERIAS
% =============================================================================
\section{Librerías: El Toolbox Estadístico}

\begin{table}[H]
\centering
\caption{Librerías de R y su propósito}
\begin{tabular}{lp{9cm}}
\toprule
\textbf{Librería} & \textbf{Para qué sirve} \\
\midrule
\texttt{readr} & Lectura eficiente de CSV. Maneja codificaciones UTF-8 \\
\texttt{dplyr} & Manipulación de datos (filtrar, seleccionar, resumir) \\
\texttt{stringr} & Trabajo con texto: extrae números de strings como ``8GB'' \\
\texttt{MASS} & Funciones avanzadas de estadística \\
\texttt{ggplot2} & Visualización con capas (geometrías, estéticas, temas) \\
\texttt{lmtest} & Pruebas de diagnóstico: Durbin-Watson, Breusch-Pagan \\
\texttt{nortest} & Pruebas de normalidad alternativas \\
\texttt{ggfortify} & Diagnósticos automáticos con \texttt{ggplot2} \\
\texttt{gridExtra} & Organiza múltiples gráficos en una ventana \\
\texttt{car} & \texttt{vif()} para detectar multicolinealidad \\
\texttt{knitr/rmarkdown} & Generación de informes dinámicos \\
\texttt{tseries} & \texttt{jarque.bera.test} para normalidad de residuos \\
\bottomrule
\end{tabular}
\end{table}

% =============================================================================
% 3. CARGA DE DATOS
% =============================================================================
\section{Carga de Datos}

\begin{lstlisting}[caption={Lectura del dataset}]
datos <- read_csv("Desktop/Final/laptop_price.csv",
                  locale = locale(encoding = "UTF-8"))
\end{lstlisting}

\textbf{Qué hace:} Lee el CSV con 1,303 laptops. El argumento \texttt{locale(encoding = "UTF-8")} es crucial: nombres de marcas y sistemas operativos contienen caracteres especiales (tildes, etc.) que sin esta configuración saldrían como basura.

% =============================================================================
% 4. LIMPIEZA Y PREPARACIÓN
% =============================================================================
\section{Limpieza y Preparación: De Texto a Números}

Este paso es \textbf{crítico}. El dataset viene con texto como ``8GB'', ``2.3kg'', ``2.5GHz''. Hay que extraer los números.

\subsection{Extracción de variables numéricas}

\begin{lstlisting}[caption={Limpiando datos}]
datos$Ram_GB <- as.numeric(gsub("GB", "", datos$Ram))
datos$Weight_kg <- as.numeric(gsub("kg", "", datos$Weight))
datos$CPU_GHz <- as.numeric(str_extract(datos$Cpu, "\\d+\\.?\\d*(?=GHz)"))
\end{lstlisting}

\textbf{Ejemplos:}
\begin{itemize}
    \item \texttt{"16GB"} $\xrightarrow{\text{gsub}}$ \texttt{"16"} $\xrightarrow{\text{as.numeric}}$ 16
    \item En \texttt{"Intel Core i7 2.50GHz"}, \texttt{str\_extract} extrae ``2.50'' usando la expresión regular: \texttt{\\d+} (uno o más dígitos), \texttt{\\.?} (punto decimal opcional), \texttt{\\d*} (dígitos extra), \texttt{(?=GHz)} (sólo si después viene GHz).
\end{itemize}

\subsection{Variables cualitativas (factores)}

\begin{lstlisting}[caption={Creando factores}]
datos$CPU_Brand <- ifelse(grepl("Intel", datos$Cpu), "Intel",
                          ifelse(grepl("AMD", datos$Cpu), "AMD", "Other"))
datos$CPU_Brand <- factor(datos$CPU_Brand)
datos$GPU_Brand <- factor(datos$GPU_Brand)
datos$TypeName <- factor(datos$TypeName)
datos$OpSys <- factor(datos$OpSys)
\end{lstlisting}

\textbf{¿Por qué \texttt{factor}?} R no puede meter texto directo en una regresión. Al convertir a \texttt{factor}, R crea automáticamente \textbf{variables dummy} (0/1) para cada categoría excepto la de referencia.

\subsection{Filtrado de datos faltantes}

\begin{lstlisting}
datos <- datos[complete.cases(
  datos$Price_euros, datos$Ram_GB, datos$CPU_GHz,
  datos$Inches, datos$TypeName, datos$OpSys), ]
\end{lstlisting}

\textbf{Resultado:} 1,303 observaciones completas. No se perdieron datos.

% =============================================================================
% 5. ESTADISTICA DESCRIPTIVA
% =============================================================================
\section{Estadística Descriptiva}

\begin{table}[H]
\centering
\caption{Resumen descriptivo de variables numéricas ($N=1{,}303$)}
\begin{tabular}{lrrrrrr}
\toprule
\textbf{Variable} & \textbf{Media} & \textbf{Desv.Std} & \textbf{Mín} & \textbf{Q1} & \textbf{Mediana} & \textbf{Q3} & \textbf{Máx} \\
\midrule
Price\_euros & 1{,}123.69 & 699.01 & 174 & 599 & 977 & 1{,}488 & 6{,}099 \\
Ram\_GB & 8.38 & 5.08 & 2 & 4 & 8 & 8 & 64 \\
CPU\_GHz & 2.30 & 0.51 & 0.90 & 2.00 & 2.50 & 2.70 & 3.60 \\
Inches & 15.02 & 1.43 & 10.1 & 14.0 & 15.6 & 15.6 & 18.4 \\
Weight\_kg & 2.04 & 0.67 & 0.69 & 1.50 & 2.04 & 2.30 & 4.70 \\
Res\_Width & 1{,}894.78 & 494.64 & 1{,}366 & 1{,}600 & 1{,}920 & 1{,}920 & 3{,}840 \\
\bottomrule
\end{tabular}
\end{table}

\textbf{Observaciones:} El precio es muy disperso (std = 699), con laptops desde 174 EUR hasta 6{,}099 EUR (Razer Blade Pro). La RAM está concentrada en 4--8 GB.

% =============================================================================
% 6. CORRELACION DE PEARSON
% =============================================================================
\section{Correlación de Pearson}

\subsection{Matriz de correlación}

\begin{table}[H]
\centering
\caption{Matriz de correlación de Pearson}
\begin{tabular}{lrrrrrr}
\toprule
 & \textbf{Precio} & \textbf{RAM} & \textbf{CPU} & \textbf{Inches} & \textbf{Peso} & \textbf{ResW} \\
\midrule
Price\_euros & 1.0000 & 0.7430 & 0.4303 & 0.0682 & 0.2104 & 0.5565 \\
Ram\_GB & 0.7430 & 1.0000 & 0.3680 & 0.2380 & 0.3839 & 0.4331 \\
CPU\_GHz & 0.4303 & 0.3680 & 1.0000 & 0.3079 & 0.3204 & 0.1835 \\
Inches & 0.0682 & 0.2380 & 0.3079 & 1.0000 & 0.8276 & -0.0712 \\
Weight\_kg & 0.2104 & 0.3839 & 0.3204 & 0.8276 & 1.0000 & -0.0329 \\
\bottomrule
\end{tabular}
\end{table}

\subsection{Pruebas de correlación ($H_0: \rho = 0$)}

\begin{table}[H]
\centering
\caption{Pruebas de significancia de la correlación}
\begin{tabular}{lrrr}
\toprule
\textbf{Par de variables} & \textbf{r} & \textbf{p-valor} & \textbf{Interpretación} \\
\midrule
RAM vs Precio & 0.7430 & $<10^{-228}$ & Correlación fuerte y altamente significativa \\
CPU\_GHz vs Precio & 0.4303 & $<10^{-59}$ & Correlación moderada, altamente significativa \\
Inches vs Precio & 0.0682 & 0.0138 & Correlación débil, apenas significativa \\
\bottomrule
\end{tabular}
\end{table}

\textbf{Conclusión:} La RAM es la variable numérica más correlacionada con el precio. El tamaño de pantalla (Inches) tiene correlación muy débil con el precio.

% =============================================================================
% 7. MODELO RLM: RESULTADOS REALES
% =============================================================================
\section{Modelo de Regresión Lineal Múltiple: Resultados Reales}

\subsection{Bondad de ajuste}

\begin{table}[H]
\centering
\caption{Métricas del modelo}
\begin{tabular}{ll}
\toprule
\textbf{Métrica} & \textbf{Valor} \\
\midrule
$R^2$ (R-squared) & 0.706 \\
$R^2$ ajustado & 0.703 \\
F-estadístico & 193.3 \\
p-valor (F) & $<0.001$ \\
Número de observaciones & 1{,}303 \\
Grados de libertad (residuos) & 1{,}286 \\
Grados de libertad (modelo) & 16 \\
\bottomrule
\end{tabular}
\end{table}

\textbf{Interpretación:} El modelo explica el \textbf{70.6\%} de la variabilidad del precio. Es un buen ajuste. El F-estadístico con p $<$ 0.001 indica que el modelo es \textbf{globalmente significativo} (al menos una variable aporta).

\subsection{Coeficientes del modelo}

\begin{table}[H]
\centering
\caption{Coeficientes estimados (significativos en negrita)}
\begin{tabular}{lrrrr}
\toprule
\textbf{Variable} & \textbf{Coef.} & \textbf{Error Std} & \textbf{t} & \textbf{p-valor} \\
\midrule
Intercepto & 42.92 & 292.17 & 0.147 & 0.883 \\
TypeName Gaming & $-46.89$ & 53.63 & $-0.874$ & 0.382 \\
TypeName Netbook & \textbf{$-247.93$} & 88.43 & $-2.804$ & \textbf{0.005} \\
TypeName Notebook & \textbf{$-293.23$} & 42.87 & $-6.841$ & \textbf{$<$0.001} \\
TypeName Ultrabook & \textbf{120.53} & 45.99 & 2.621 & \textbf{0.009} \\
TypeName Workstation & \textbf{650.14} & 85.03 & 7.646 & \textbf{$<$0.001} \\
OpSys Chrome OS & 263.41 & 282.90 & 0.931 & 0.352 \\
OpSys Linux & 231.80 & 279.93 & 0.828 & 0.408 \\
OpSys Mac OS X & 401.45 & 306.28 & 1.311 & 0.190 \\
OpSys No OS & 179.50 & 279.59 & 0.642 & 0.521 \\
OpSys Windows 10 & 420.91 & 274.98 & 1.531 & 0.126 \\
OpSys Windows 10 S & 502.64 & 306.60 & 1.639 & 0.101 \\
OpSys Windows 7 & \textbf{827.97} & 280.87 & 2.948 & \textbf{0.003} \\
OpSys macOS & \textbf{617.62} & 295.90 & 2.087 & \textbf{0.037} \\
\textbf{Ram\_GB} & \textbf{81.08} & 2.58 & 31.414 & \textbf{$<$0.001} \\
\textbf{CPU\_GHz} & \textbf{208.31} & 24.49 & 8.505 & \textbf{$<$0.001} \\
\textbf{Inches} & \textbf{$-23.09$} & 11.15 & $-2.071$ & \textbf{0.039} \\
\bottomrule
\end{tabular}
\end{table}

\textbf{Categorías de referencia:}
\begin{itemize}
    \item TypeName: \textbf{``2 in 1 Convertible''} (la primera alfabéticamente)
    \item OpSys: \textbf{``Android''} (la primera alfabéticamente)
\end{itemize}

\subsection{Interpretación de coeficientes}

\subsubsection{Variables numéricas}

\begin{itemize}
    \item \textbf{RAM ($\beta = 81.08$):} Por cada GB adicional de RAM, el precio aumenta en promedio \textbf{81.08 euros}, manteniendo constantes CPU, pulgadas, tipo y sistema operativo.
    
    \item \textbf{CPU\_GHz ($\beta = 208.31$):} Por cada GHz adicional de frecuencia del procesador, el precio sube en promedio \textbf{208.31 euros} (ceteris paribus).
    
    \item \textbf{Inches ($\beta = -23.09$):} Cada pulgada adicional de pantalla reduce el precio en promedio \textbf{23.09 euros}. Esto es contra-intuitivo pero tiene explicación: las laptops más grandes tienden a ser económicas (Chromebooks, notebooks básicas), mientras que las ultrabooks premium son más pequeñas y caras.
\end{itemize}

\subsubsection{Variables cualitativas (dummy)}

\textbf{Con respecto a ``2 in 1 Convertible'' (referencia):}
\begin{itemize}
    \item Un \textbf{Netbook} cuesta 247.93 euros \textbf{menos} que un ``2 in 1 Convertible'' equivalente.
    \item Un \textbf{Notebook} cuesta 293.23 euros \textbf{menos}.
    \item Un \textbf{Ultrabook} cuesta 120.53 euros \textbf{más}.
    \item Una \textbf{Workstation} cuesta 650.14 euros \textbf{más}.
    \item Un \textbf{Gaming} no es significativamente diferente del ``2 in 1 Convertible'' (p = 0.382).
\end{itemize}

\textbf{Con respecto a ``Android'' (referencia de sistema operativo):}
\begin{itemize}
    \item Windows 7 y macOS son significativamente más caros que Android (p = 0.003 y p = 0.037 respectivamente).
    \item Windows 10 no es significativamente diferente de Android (p = 0.126).
\end{itemize}

% =============================================================================
% 8. VALIDACION DE SUPUESTOS
% =============================================================================
\section{Validación de Supuestos: Resultados Reales}

\subsection{1. Linealidad}

Se verifica con gráficos de residuos vs cada variable independiente. Si la línea suavizada se mantiene cerca del cero horizontal, el supuesto se cumple.

\textbf{Veredicto:} La linealidad parece razonable en los gráficos exploratorios.

\subsection{2. Normalidad de residuos}

\begin{table}[H]
\centering
\caption{Pruebas de normalidad de residuos}
\begin{tabular}{lrrr}
\toprule
\textbf{Prueba} & \textbf{Estadístico} & \textbf{p-valor} & \textbf{Conclusión} \\
\midrule
Shapiro-Wilk & $W = 0.9238$ & $<0.001$ & \textcolor{rojoacento}{NO Normalidad} \\
Jarque-Bera & $JB = 2411.46$ & $<0.001$ & \textcolor{rojoacento}{NO Normalidad} \\
\bottomrule
\end{tabular}
\end{table}

\textbf{Interpretación:} Los residuos \textbf{no siguen una distribución normal}. Esto se debe a que el precio tiene colas largas (asimetría positiva: pocos laptops muy caros como Razer a 6,099 EUR).

\textbf{¿Qué hacer?} 
\begin{enumerate}
    \item Transformar la variable dependiente con logaritmo: $Y' = \ln(Y)$.
    \item Usar errores estándar robustos (White/HC3) que no dependen de normalidad.
    \item Aceptar que con $N = 1{,}303$, el Teorema Central del Límite aplica y los p-valores de los coeficientes siguen siendo aproximadamente válidos.
\end{enumerate}

\subsection{3. Independencia de residuos}

\begin{table}[H]
\centering
\caption{Prueba de independencia}
\begin{tabular}{lrr}
\toprule
\textbf{Prueba} & \textbf{Estadístico} & \textbf{Conclusión} \\
\midrule
Durbin-Watson & $DW = 2.027$ & \textcolor{verdeok}{Independencia cumplida} \\
\bottomrule
\end{tabular}
\end{table}

\textbf{Interpretación:} $DW \approx 2$ indica ausencia de autocorrelación. Los residuos son independientes.

\subsection{4. Homocedasticidad}

\begin{table}[H]
\centering
\caption{Prueba de homocedasticidad}
\begin{tabular}{lrrr}
\toprule
\textbf{Prueba} & \textbf{Estadístico} & \textbf{p-valor} & \textbf{Conclusión} \\
\midrule
Breusch-Pagan & $LM = 201.21$ & $<0.001$ & \textcolor{rojoacento}{HETEROCEDASTICIDAD} \\
\bottomrule
\end{tabular}
\end{table}

\textbf{Interpretación:} Existe heterocedasticidad. La varianza del error \textbf{no es constante} para todos los niveles de las variables. Esto es común en datos de precios: las laptops caras tienen mayor variabilidad de precio que las baratas.

\textbf{¿Qué hacer?}
\begin{enumerate}
    \item Transformar $Y$ con logaritmo.
    \item Usar errores estándar robustos (HC3).
    \item No ignorarlo: los p-valores actuales \textbf{podrían estar ligeramente sesgados}.
\end{enumerate}

\subsection{5. Multicolinealidad (VIF)}

\begin{table}[H]
\centering
\caption{Factor de Inflación de Varianza (VIF)}
\begin{tabular}{lrr}
\toprule
\textbf{Variable} & \textbf{VIF} & \textbf{Alerta} \\
\midrule
Ram\_GB & 1.54 & \textcolor{verdeok}{OK} \\
CPU\_GHz & 1.38 & \textcolor{verdeok}{OK} \\
Inches & 2.27 & \textcolor{verdeok}{OK} \\
TypeName Gaming & 3.42 & \textcolor{verdeok}{OK} \\
TypeName Netbook & 1.32 & \textcolor{verdeok}{OK} \\
TypeName Notebook & 4.06 & \textcolor{verdeok}{OK} \\
TypeName Ultrabook & 2.42 & \textcolor{verdeok}{OK} \\
TypeName Workstation & 1.41 & \textcolor{verdeok}{OK} \\
OpSys Chrome OS & 14.57 & \textcolor{rojoacento}{\textbf{ALTO}} \\
OpSys Linux & 31.85 & \textcolor{rojoacento}{\textbf{MUY ALTO}} \\
OpSys Mac OS X & 5.13 & \textcolor{verdeok}{Moderado} \\
OpSys No OS & 33.71 & \textcolor{rojoacento}{\textbf{MUY ALTO}} \\
OpSys Windows 10 & 98.91 & \textcolor{rojoacento}{\textbf{SEVERO}} \\
OpSys Windows 10 S & 5.14 & \textcolor{verdeok}{Moderado} \\
OpSys Windows 7 & 23.59 & \textcolor{rojoacento}{\textbf{MUY ALTO}} \\
OpSys macOS & 7.76 & \textcolor{rojoacento}{\textbf{ALTO}} \\
\bottomrule
\end{tabular}
\end{table}

\textbf{¿Por qué pasa esto?} La variable OpSys tiene \textbf{demasiadas categorías} con pocas observaciones en algunas (ej. Android solo 1 laptop, Chrome OS pocas). Esto genera colinealidad artificial: si sabemos que NO es Windows 10, es más probable que sea Linux, etc. Los VIF altos significan que los errores estándar de los coeficientes de OpSys están \textbf{inflados} y menos confiables.

\textbf{¿Qué hacer?}
\begin{enumerate}
    \item \textbf{Recategorizar} OpSys en 3-4 grupos: Windows, macOS, Linux/No OS, Chrome OS.
    \item Eliminar categorías con muy pocas observaciones.
    \item Usar errores estándar robustos para mitigar el efecto.
\end{enumerate}

% =============================================================================
% 9. INTERVALOS DE CONFIANZA
% =============================================================================
\section{Intervalos de Confianza al 95\%}

\begin{table}[H]
\centering
\caption{Intervalos de confianza para coeficientes significativos}
\begin{tabular}{lrr}
\toprule
\textbf{Variable} & \textbf{Límite inferior} & \textbf{Límite superior} \\
\midrule
Netbook & $-421.40$ & $-74.45$ \\
Notebook & $-377.33$ & $-209.14$ \\
Ultrabook & 30.31 & 210.75 \\
Workstation & 483.32 & 816.96 \\
Windows 7 & 276.96 & 1{,}378.97 \\
macOS & 37.13 & 1{,}198.12 \\
Ram\_GB & 76.01 & 86.14 \\
CPU\_GHz & 160.26 & 256.36 \\
Inches & $-44.96$ & $-1.21$ \\
\bottomrule
\end{tabular}
\end{table}

\textbf{Nota:} Si un intervalo incluye el cero, el coeficiente NO es significativo (consistente con p-valor $>$ 0.05).

% =============================================================================
% 10. PREDICCION
% =============================================================================
\section{Predicción con el Modelo}

Se predice el precio para una laptop con estas características:
\begin{itemize}
    \item RAM = 16 GB
    \item CPU = 2.8 GHz
    \item Pantalla = 15.6 pulgadas
    \item Tipo = Gaming
    \item Sistema operativo = Windows 10
\end{itemize}

\begin{table}[H]
\centering
\caption{Predicción}
\begin{tabular}{lr}
\toprule
\textbf{Concepto} & \textbf{Valor (EUR)} \\
\midrule
Predicción puntual & 1{,}937.24 \\
Límite inferior (95\%) & 1{,}881.05 \\
Límite superior (95\%) & 1{,}993.43 \\
\bottomrule
\end{tabular}
\end{table}

\textbf{Interpretación:} Se predice que una laptop con esas características costaría aproximadamente \textbf{1,937 EUR}, con un intervalo de confianza al 95\% entre 1,881 y 1,993 EUR.

% =============================================================================
% 11. PROBLEMAS Y RECOMENDACIONES
% =============================================================================
\section{Problemas Detectados y Recomendaciones}

\subsection{Problemas}

\begin{enumerate}
    \item \textbf{No normalidad de residuos:} Los residuos tienen colas pesadas (kurtosis = 9.26, asimetría = 1.15).
    
    \item \textbf{Heterocedasticidad:} La varianza del error no es constante (Breusch-Pagan p $<$ 0.001).
    
    \item \textbf{Multicolinealidad severa en OpSys:} VIF de Windows 10 = 98.91 (severo). Esto ocurre porque OpSys tiene demasiadas categorías con distribución desigual.
\end{enumerate}

\subsection{Recomendaciones para mejorar el modelo}

\begin{enumerate}
    \item \textbf{Transformar el precio con logaritmo:} $Y' = \ln(\text{Price\_euros})$. Esto suele corregir la asimetría y estabilizar la varianza, resolviendo normalidad y heterocedasticidad simultáneamente.
    
    \item \textbf{Recategorizar OpSys:} Agrupar en 4 categorías: Windows, macOS, Linux/No OS, Chrome OS/Android. Esto reducirá drásticamente los VIF.
    
    \item \textbf{Usar errores estándar robustos:} Con \texttt{lmtest + sandwich} o \texttt{estimatr::lm\_robust} para que los p-valores sean válidos a pesar de heterocedasticidad.
    
    \item \textbf{Considerar incluir más variables:} Marca (Company), GPU\_Brand, peso, resolución de pantalla.
\end{enumerate}

% =============================================================================
% 12. CONCLUSIONES DEL ANÁLISIS
% =============================================================================
\section{Conclusiones del Análisis}

\begin{enumerate}
    \item \textbf{La RAM es el factor más importante:} Con un coeficiente de 81.08 EUR por GB y p-valor $< 10^{-228}$, es la variable numérica con mayor impacto sobre el precio.
    
    \item \textbf{El tipo de laptop importa mucho:} Workstations cuestan 650 EUR más que ``2 in 1 Convertible'', Notebooks 293 EUR menos. El tipo de equipo captura diferencias de mercado importantes.
    
    \item \textbf{El sistema operativo tiene problemas de multicolinealidad:} Los VIF altos sugieren que las categorías de OpSys están demasiado fragmentadas. Una recategorización mejoraría el modelo.
    
    \item \textbf{Los supuestos no se cumplen perfectamente:} La no normalidad y la heterocedasticidad indican que el modelo podría mejorarse con una transformación logarítmica del precio.
    
    \item \textbf{El $R^2 = 0.706$ es bueno pero hay margen:} Con 1,303 observaciones, un 70\% de varianza explicada es sólido, pero variables como marca, GPU y almacenamiento podrían elevarlo más.
\end{enumerate}

% =============================================================================
% 13. CHECKLIST PARA LA PRESENTACIÓN
% =============================================================================
\section{Checklist para la Presentación}

Antes de exponer, domina estas respuestas:

\begin{enumerate}
    \item ¿De dónde salen los datos? ¿Cuántas observaciones tienes? \textbf{R:} Kaggle, 1,303 laptops.
    
    \item ¿Qué hace \texttt{str\_extract} y por qué fue necesario? \textbf{R:} Extrae números de texto usando expresiones regulares; porque las columnas vienen como ``8GB'', ``2.5GHz''.
    
    \item ¿Por qué convertiste \texttt{TypeName} y \texttt{OpSys} a \texttt{factor}? \textbf{R:} Para que R cree variables dummy automáticamente y las incluya en la regresión.
    
    \item ¿Cuál es la interpretación del coeficiente de RAM? \textbf{R:} Por cada GB adicional, el precio aumenta 81.08 EUR en promedio, ceteris paribus.
    
    \item ¿Qué significa que un coeficiente tenga p-valor $<$ 0.05? \textbf{R:} Que es estadísticamente significativo; es improbable que ese efecto sea cero en la población.
    
    \item ¿Qué es el $R^2$ ajustado y por qué es más honesto? \textbf{R:} Penaliza agregar variables inútiles; sube solo si la nueva variable realmente aporta.
    
    \item ¿Cuáles son los 5 supuestos de la RLM? \textbf{R:} Linealidad, normalidad de residuos, independencia, homocedasticidad, no multicolinealidad.
    
    \item ¿Qué muestra el Q-Q plot y qué significa si los puntos se desvían? \textbf{R:} Compara cuantiles de tus residuos vs. una normal. Si se desvían, los residuos no son normales.
    
    \item ¿Qué es la heterocedasticidad y cómo se corrige? \textbf{R:} Varianza del error no constante. Se corrige con transformación log o errores estándar robustos.
    
    \item ¿Qué es el VIF y cuál es el umbral de alarma? \textbf{R:} Factor de inflación de varianza. VIF $>$ 10 es severo; entre 5 y 10 es moderado.
    
    \item ¿Por qué los VIF de OpSys están tan altos? \textbf{R:} Porque hay demasiadas categorías con pocas observaciones, generando colinealidad artificial.
    
    \item ¿Cómo se interpreta el coeficiente de una variable dummy? \textbf{R:} Es la diferencia de precio promedio respecto a la categoría de referencia, manteniendo todo lo demás constante.
    
    \item Si te piden predecir un nuevo equipo, ¿qué precaución tienes? \textbf{R:} Las variables cualitativas deben ser \texttt{factor} con los mismos niveles que el dataset de entrenamiento.
\end{enumerate}

% =============================================================================
% FIN
% =============================================================================
\end{document}
